\documentclass[tikz,border=10pt]{standalone}
\usepackage[utf8]{inputenc}
\usepackage[T1]{fontenc}
\usepackage[lithuanian]{babel}
\usetikzlibrary{shapes.geometric, arrows.meta, positioning, shadows, fit, backgrounds}

\hyphenpenalty=10000
\exhyphenpenalty=10000

\definecolor{mainblue}{RGB}{41, 98, 255}
\definecolor{cat1}{RGB}{66, 165, 245}
\definecolor{cat2}{RGB}{102, 187, 106}
\definecolor{cat3}{RGB}{255, 167, 38}
\definecolor{cat4}{RGB}{171, 71, 188}
\definecolor{cat5}{RGB}{239, 83, 80}
\definecolor{cat6}{RGB}{38, 166, 154}
\definecolor{cat7}{RGB}{121, 134, 203}
\definecolor{cat8}{RGB}{255, 112, 67}
\definecolor{cat9}{RGB}{158, 157, 36}
\definecolor{cat10}{RGB}{0, 151, 167}
\definecolor{cat11}{RGB}{216, 27, 96}
\definecolor{cat12}{RGB}{109, 76, 65}

\tikzstyle{maintitle} = [rectangle, rounded corners=8pt, minimum width=5cm, minimum height=1.2cm,
    text centered, draw=mainblue, fill=mainblue, text=white, font=\bfseries\Large, drop shadow]
\tikzstyle{category} = [rectangle, rounded corners=5pt, minimum width=3.6cm, minimum height=0.9cm,
    text centered, font=\bfseries\scriptsize, drop shadow, text width=3.4cm]
\tikzstyle{arrow} = [thick, ->, >=Stealth]

\begin{document}
\begin{tikzpicture}[node distance=0.8cm]

% Pagrindinis pavadinimas
\node (main) [maintitle, text width=14cm] {Mokslinių duomenų klasifikacija pagal generavimo pobūdį};

% 1 eilutės jungiamoji horizontali linija
\coordinate (junc1) at ([yshift=-1cm]main.south);
\draw [mainblue, thick] ([xshift=-9.5cm]junc1) -- ([xshift=9.5cm]junc1);
\draw [mainblue, thick] (main.south) -- (junc1);

% 1 eilutė – 6 kategorijos
\node (c1) [category, below=2.2cm of main, xshift=-9.5cm, draw=cat1, fill=cat1!15] {Eksperimentiniai duomenys};
\node (c2) [category, below=2.2cm of main, xshift=-5.7cm, draw=cat2, fill=cat2!15] {Stebėjimų duomenys};
\node (c3) [category, below=2.2cm of main, xshift=-1.9cm, draw=cat3, fill=cat3!15] {Apklausų duomenys};
\node (c4) [category, below=2.2cm of main, xshift=1.9cm, draw=cat4, fill=cat4!15] {Skaičiavimų / modeliavimo duomenys};
\node (c5) [category, below=2.2cm of main, xshift=5.7cm, draw=cat5, fill=cat5!15] {Tekstiniai ir dokumentiniai duomenys};
\node (c6) [category, below=2.2cm of main, xshift=9.5cm, draw=cat6, fill=cat6!15] {Vaizdiniai duomenys};

% Rodyklės į 1 eilutę
\draw [arrow, cat1] ([xshift=-9.5cm]junc1) -- (c1.north);
\draw [arrow, cat2] ([xshift=-5.7cm]junc1) -- (c2.north);
\draw [arrow, cat3] ([xshift=-1.9cm]junc1) -- (c3.north);
\draw [arrow, cat4] ([xshift=1.9cm]junc1) -- (c4.north);
\draw [arrow, cat5] ([xshift=5.7cm]junc1) -- (c5.north);
\draw [arrow, cat6] ([xshift=9.5cm]junc1) -- (c6.north);

% 2 eilutės jungiamoji horizontali linija (po 1 eilutės apačia)
\coordinate (junc2) at ([yshift=-4cm]main.south);
\draw [mainblue, thick] ([xshift=-9.5cm]junc2) -- ([xshift=9.5cm]junc2);
\draw [mainblue, thick] (junc1) -- (junc2);

% 2 eilutė – 6 kategorijos
\node (c7) [category, below=5.2cm of main, xshift=-9.5cm, draw=cat7, fill=cat7!15] {Erdviniai duomenys};
\node (c8) [category, below=5.2cm of main, xshift=-5.7cm, draw=cat8, fill=cat8!15] {Audio ir video duomenys};
\node (c9) [category, below=5.2cm of main, xshift=-1.9cm, draw=cat9, fill=cat9!15] {Išvestiniai (antriniai) duomenys};
\node (c10) [category, below=5.2cm of main, xshift=1.9cm, draw=cat10, fill=cat10!15] {Programinės įrangos ir kodo duomenys};
\node (c11) [category, below=5.2cm of main, xshift=5.7cm, draw=cat11, fill=cat11!15] {Interaktyvūs ir 3D ištekliai};
\node (c12) [category, below=5.2cm of main, xshift=9.5cm, draw=cat12, fill=cat12!15] {Skaitmenizuoti archyviniai šaltiniai};

% Rodyklės į 2 eilutę
\draw [arrow, cat7] ([xshift=-9.5cm]junc2) -- (c7.north);
\draw [arrow, cat8] ([xshift=-5.7cm]junc2) -- (c8.north);
\draw [arrow, cat9] ([xshift=-1.9cm]junc2) -- (c9.north);
\draw [arrow, cat10] ([xshift=1.9cm]junc2) -- (c10.north);
\draw [arrow, cat11] ([xshift=5.7cm]junc2) -- (c11.north);
\draw [arrow, cat12] ([xshift=9.5cm]junc2) -- (c12.north);

\end{tikzpicture}
\end{document}
